\IfFileExists{stacks-project.cls}{%
\documentclass{stacks-project}
}{%
\documentclass{amsart}
}

% For dealing with references we use the comment environment
\usepackage{verbatim}
\newenvironment{reference}{\comment}{\endcomment}
%\newenvironment{reference}{}{}
\newenvironment{slogan}{\comment}{\endcomment}
\newenvironment{history}{\comment}{\endcomment}

% For commutative diagrams we use Xy-pic
\usepackage[all]{xy}

% We use 2cell for 2-commutative diagrams.
\xyoption{2cell}
\UseAllTwocells

% We use multicol for the list of chapters between chapters
\usepackage{multicol}

% This is generally recommended for better output
\usepackage{lmodern}
\usepackage[T1]{fontenc}

% For cross-file-references
\usepackage{xr-hyper}

% Package for hypertext links:
\usepackage{hyperref}

% For any local file, say "hello.tex" you want to link to please
% use \externaldocument[hello-]{hello}
\externaldocument[introduction-]{introduction}
\externaldocument[conventions-]{conventions}
\externaldocument[sets-]{sets}
\externaldocument[categories-]{categories}
\externaldocument[topology-]{topology}
\externaldocument[sheaves-]{sheaves}
\externaldocument[sites-]{sites}
\externaldocument[stacks-]{stacks}
\externaldocument[fields-]{fields}
\externaldocument[algebra-]{algebra}
\externaldocument[brauer-]{brauer}
\externaldocument[homology-]{homology}
\externaldocument[derived-]{derived}
\externaldocument[simplicial-]{simplicial}
\externaldocument[more-algebra-]{more-algebra}
\externaldocument[smoothing-]{smoothing}
\externaldocument[modules-]{modules}
\externaldocument[sites-modules-]{sites-modules}
\externaldocument[injectives-]{injectives}
\externaldocument[cohomology-]{cohomology}
\externaldocument[sites-cohomology-]{sites-cohomology}
\externaldocument[dga-]{dga}
\externaldocument[dpa-]{dpa}
\externaldocument[sdga-]{sdga}
\externaldocument[hypercovering-]{hypercovering}
\externaldocument[schemes-]{schemes}
\externaldocument[constructions-]{constructions}
\externaldocument[properties-]{properties}
\externaldocument[morphisms-]{morphisms}
\externaldocument[coherent-]{coherent}
\externaldocument[divisors-]{divisors}
\externaldocument[limits-]{limits}
\externaldocument[varieties-]{varieties}
\externaldocument[topologies-]{topologies}
\externaldocument[descent-]{descent}
\externaldocument[perfect-]{perfect}
\externaldocument[more-morphisms-]{more-morphisms}
\externaldocument[flat-]{flat}
\externaldocument[groupoids-]{groupoids}
\externaldocument[more-groupoids-]{more-groupoids}
\externaldocument[etale-]{etale}
\externaldocument[chow-]{chow}
\externaldocument[intersection-]{intersection}
\externaldocument[pic-]{pic}
\externaldocument[weil-]{weil}
\externaldocument[adequate-]{adequate}
\externaldocument[dualizing-]{dualizing}
\externaldocument[duality-]{duality}
\externaldocument[discriminant-]{discriminant}
\externaldocument[derham-]{derham}
\externaldocument[local-cohomology-]{local-cohomology}
\externaldocument[algebraization-]{algebraization}
\externaldocument[curves-]{curves}
\externaldocument[resolve-]{resolve}
\externaldocument[models-]{models}
\externaldocument[functors-]{functors}
\externaldocument[equiv-]{equiv}
\externaldocument[pione-]{pione}
\externaldocument[etale-cohomology-]{etale-cohomology}
\externaldocument[proetale-]{proetale}
\externaldocument[relative-cycles-]{relative-cycles}
\externaldocument[more-etale-]{more-etale}
\externaldocument[trace-]{trace}
\externaldocument[crystalline-]{crystalline}
\externaldocument[spaces-]{spaces}
\externaldocument[spaces-properties-]{spaces-properties}
\externaldocument[spaces-morphisms-]{spaces-morphisms}
\externaldocument[decent-spaces-]{decent-spaces}
\externaldocument[spaces-cohomology-]{spaces-cohomology}
\externaldocument[spaces-limits-]{spaces-limits}
\externaldocument[spaces-divisors-]{spaces-divisors}
\externaldocument[spaces-over-fields-]{spaces-over-fields}
\externaldocument[spaces-topologies-]{spaces-topologies}
\externaldocument[spaces-descent-]{spaces-descent}
\externaldocument[spaces-perfect-]{spaces-perfect}
\externaldocument[spaces-more-morphisms-]{spaces-more-morphisms}
\externaldocument[spaces-flat-]{spaces-flat}
\externaldocument[spaces-groupoids-]{spaces-groupoids}
\externaldocument[spaces-more-groupoids-]{spaces-more-groupoids}
\externaldocument[bootstrap-]{bootstrap}
\externaldocument[spaces-pushouts-]{spaces-pushouts}
\externaldocument[spaces-chow-]{spaces-chow}
\externaldocument[groupoids-quotients-]{groupoids-quotients}
\externaldocument[spaces-more-cohomology-]{spaces-more-cohomology}
\externaldocument[spaces-simplicial-]{spaces-simplicial}
\externaldocument[spaces-duality-]{spaces-duality}
\externaldocument[formal-spaces-]{formal-spaces}
\externaldocument[restricted-]{restricted}
\externaldocument[spaces-resolve-]{spaces-resolve}
\externaldocument[formal-defos-]{formal-defos}
\externaldocument[defos-]{defos}
\externaldocument[cotangent-]{cotangent}
\externaldocument[examples-defos-]{examples-defos}
\externaldocument[algebraic-]{algebraic}
\externaldocument[examples-stacks-]{examples-stacks}
\externaldocument[stacks-sheaves-]{stacks-sheaves}
\externaldocument[criteria-]{criteria}
\externaldocument[artin-]{artin}
\externaldocument[quot-]{quot}
\externaldocument[stacks-properties-]{stacks-properties}
\externaldocument[stacks-morphisms-]{stacks-morphisms}
\externaldocument[stacks-limits-]{stacks-limits}
\externaldocument[stacks-cohomology-]{stacks-cohomology}
\externaldocument[stacks-perfect-]{stacks-perfect}
\externaldocument[stacks-introduction-]{stacks-introduction}
\externaldocument[stacks-more-morphisms-]{stacks-more-morphisms}
\externaldocument[stacks-geometry-]{stacks-geometry}
\externaldocument[moduli-]{moduli}
\externaldocument[moduli-curves-]{moduli-curves}
\externaldocument[examples-]{examples}
\externaldocument[exercises-]{exercises}
\externaldocument[guide-]{guide}
\externaldocument[desirables-]{desirables}
\externaldocument[coding-]{coding}
\externaldocument[obsolete-]{obsolete}
\externaldocument[fdl-]{fdl}
\externaldocument[index-]{index}

% Theorem environments.
%
\theoremstyle{plain}
\newtheorem{theorem}[subsection]{Theorem}
\newtheorem{proposition}[subsection]{Proposition}
\newtheorem{lemma}[subsection]{Lemma}

\theoremstyle{definition}
\newtheorem{definition}[subsection]{Definition}
\newtheorem{example}[subsection]{Example}
\newtheorem{exercise}[subsection]{Exercise}
\newtheorem{situation}[subsection]{Situation}

\theoremstyle{remark}
\newtheorem{remark}[subsection]{Remark}
\newtheorem{remarks}[subsection]{Remarks}

\numberwithin{equation}{subsection}

% Macros
%
\def\lim{\mathop{\mathrm{lim}}\nolimits}
\def\colim{\mathop{\mathrm{colim}}\nolimits}
\def\Spec{\mathop{\mathrm{Spec}}}
\def\Hom{\mathop{\mathrm{Hom}}\nolimits}
\def\Ext{\mathop{\mathrm{Ext}}\nolimits}
\def\SheafHom{\mathop{\mathcal{H}\!\mathit{om}}\nolimits}
\def\SheafExt{\mathop{\mathcal{E}\!\mathit{xt}}\nolimits}
\def\Sch{\mathit{Sch}}
\def\Mor{\mathop{\mathrm{Mor}}\nolimits}
\def\Ob{\mathop{\mathrm{Ob}}\nolimits}
\def\Sh{\mathop{\mathit{Sh}}\nolimits}
\def\NL{\mathop{N\!L}\nolimits}
\def\CH{\mathop{\mathrm{CH}}\nolimits}
\def\proetale{{pro\text{-}\acute{e}tale}}
\def\etale{{\acute{e}tale}}
\def\QCoh{\mathit{QCoh}}
\def\Ker{\mathop{\mathrm{Ker}}}
\def\Im{\mathop{\mathrm{Im}}}
\def\Coker{\mathop{\mathrm{Coker}}}
\def\Coim{\mathop{\mathrm{Coim}}}

% Boxtimes
%
\DeclareMathSymbol{\boxtimes}{\mathbin}{AMSa}{"02}

%
% Macros for moduli stacks/spaces
%
\def\QCohstack{\mathcal{QC}\!\mathit{oh}}
\def\Cohstack{\mathcal{C}\!\mathit{oh}}
\def\Spacesstack{\mathcal{S}\!\mathit{paces}}
\def\Quotfunctor{\mathrm{Quot}}
\def\Hilbfunctor{\mathrm{Hilb}}
\def\Curvesstack{\mathcal{C}\!\mathit{urves}}
\def\Polarizedstack{\mathcal{P}\!\mathit{olarized}}
\def\Complexesstack{\mathcal{C}\!\mathit{omplexes}}
% \Pic is the operator that assigns to X its picard group, usage \Pic(X)
% \Picardstack_{X/B} denotes the Picard stack of X over B
% \Picardfunctor_{X/B} denotes the Picard functor of X over B
\def\Pic{\mathop{\mathrm{Pic}}\nolimits}
\def\Picardstack{\mathcal{P}\!\mathit{ic}}
\def\Picardfunctor{\mathrm{Pic}}
\def\Deformationcategory{\mathcal{D}\!\mathit{ef}}
\usepackage{cleveref}
\crefname{originaldefinition}{original definition}{original definitions}
\Crefname{originaldefinition}{Original definition}{Original definitions}
\makeatletter
\newlabel{flat-definition-h-covering}{{0ETS}{source}{}{sourcepointer-1}{}}
\newlabel{flat-definition-h-covering@cref}{{[originaldefinition][0][]0ETS}{[1][0][]source}{}{}{}}
\makeatother
\crefname{originallemma}{original lemma}{original lemmas}
\Crefname{originallemma}{Original lemma}{Original lemmas}
\makeatletter
\newlabel{lemma-relative-cycle-functoriality}{{0H51}{source}{}{sourcepointer-2}{}}
\newlabel{lemma-relative-cycle-functoriality@cref}{{[originallemma][0][]0H51}{[1][0][]source}{}{}{}}
\makeatother
\crefname{originallemma}{original lemma}{original lemmas}
\Crefname{originallemma}{Original lemma}{Original lemmas}
\makeatletter
\newlabel{flat-lemma-h}{{0ETY}{source}{}{sourcepointer-3}{}}
\newlabel{flat-lemma-h@cref}{{[originallemma][0][]0ETY}{[1][0][]source}{}{}{}}
\makeatother
\crefname{originallemma}{original lemma}{original lemmas}
\Crefname{originallemma}{Original lemma}{Original lemmas}
\makeatletter
\newlabel{topologies-lemma-refine-qcqs-V}{{0ETN}{source}{}{sourcepointer-4}{}}
\newlabel{topologies-lemma-refine-qcqs-V@cref}{{[originallemma][0][]0ETN}{[1][0][]source}{}{}{}}
\makeatother
\crefname{originallemma}{original lemma}{original lemmas}
\Crefname{originallemma}{Original lemma}{Original lemmas}
\makeatletter
\newlabel{properties-lemma-locally-Noetherian-specialization-dvr}{{054F}{source}{}{sourcepointer-5}{}}
\newlabel{properties-lemma-locally-Noetherian-specialization-dvr@cref}{{[originallemma][0][]054F}{[1][0][]source}{}{}{}}
\makeatother
\crefname{originallemma}{original lemma}{original lemmas}
\Crefname{originallemma}{Original lemma}{Original lemmas}
\makeatletter
\newlabel{lemma-specialization-extension}{{0H4K}{source}{}{sourcepointer-6}{}}
\newlabel{lemma-specialization-extension@cref}{{[originallemma][0][]0H4K}{[1][0][]source}{}{}{}}
\makeatother
\crefname{originalsection}{original section}{original sections}
\Crefname{originalsection}{Original section}{Original sections}
\makeatletter
\newlabel{section-relative-fields}{{0H4E}{source}{}{sourcepointer-7}{}}
\newlabel{section-relative-fields@cref}{{[originalsection][0][]0H4E}{[1][0][]source}{}{}{}}
\makeatother
\title{The relative cycle property satisfies h-descent}
\author{The Stacks Project Authors}
\date{stacks-project:relative-cycles.tex:local-snapshot-20261012; GNU FDL 1.2 or later}
\begin{document}
\maketitle
\setcounter{subsection}{1}
\noindent\textit{Editorial note.} All delivered mathematical text is the author's own, copied from the pinned Stacks Project chapter file relative-cycles.tex; no statement or proof step was rewritten, reordered or removed. The counted span is the author environment and its own proof environment. References inside the span are handed to the assembler as navigation stubs that carry the canonical tag of the cited tag only; whether each of them is load-bearing is recorded on the stub itself.\par\medskip
\noindent\textit{Allowed standard prerequisites (editorial).}\begin{itemize}
\item Functoriality of relative cycles under pullback (relative-cycles chapter, Lemma lemma-relative-cycle-functoriality, canonical tag 0H51). cited at proof lines 990-991. Base changes of relative cycles are relative cycles, proving the forward direction.
\item The h covering notion and stability of h coverings under base change (More on Flatness, Definition flat-definition-h-covering, canonical tag 0ETS, and Lemma flat-lemma-h, canonical tag 0ETY). cited at statement lines 984-985 and proof lines 996-997. After replacing S by the valuation ring the pulled back family \{S\_i x\_S S' to S'\} is again an h covering.
\item Refinement of qcqs V coverings (Topologies, Lemma topologies-lemma-refine-qcqs-V, canonical tag 0ETN). cited at proof lines 1002-1005. Since an h covering is a V covering, some member S\_i carries points s' specializing to s lying over the generic and closed points of the valuation ring.
\item Realizing specializations by valuation rings (Properties chapter, Lemma properties-lemma-locally-Noetherian-specialization-dvr, canonical tag 054F). cited at proof lines 1006-1008. Provides h : S' to S\_i from a valuation ring mapping its generic and closed points to s' and s.
\item Compatibility of specialization maps with base change along extensions of discrete valuation rings (relative-cycles chapter, Lemma lemma-specialization-extension, canonical tag 0H4K). cited at proof lines 1012-1015. The composite g = g\_i composed with h is induced by a DVR extension, so the specialization identity transports from the base change to the original family.
\item Injectivity of base change for cycles over fields (relative-cycles chapter, Section section-relative-fields, canonical tag 0H4E). cited at proof lines 1017-1018. The transported identity on the valuation ring base change descends because the closed-fibre base change map is injective.
\end{itemize}\medskip
\ifdefined\hypertarget\hypertarget{sourcepointer-1}{}\fi
\noindent\textit{Source reference.} Original definition 0ETS: sources/stacks-chow/tags.txt, source lines 18527--18527. Cited at statement lines 984-985: the covering class in the statement.\par\medskip
\ifdefined\hypertarget\hypertarget{sourcepointer-2}{}\fi
\noindent\textit{Source reference.} Original lemma 0H51: sources/stacks-chow/tags.txt, source lines 21371--21371. Cited at proof lines 990-991: forward direction.\par\medskip
\ifdefined\hypertarget\hypertarget{sourcepointer-3}{}\fi
\noindent\textit{Source reference.} Original lemma 0ETY: sources/stacks-chow/tags.txt, source lines 18533--18533. Cited at proof lines 996-997: the pulled back covering is again an h covering.\par\medskip
\ifdefined\hypertarget\hypertarget{sourcepointer-4}{}\fi
\noindent\textit{Source reference.} Original lemma 0ETN: sources/stacks-chow/tags.txt, source lines 18523--18523. Cited at proof lines 1002-1005: extracts a member with the specialization over the two points of the valuation ring.\par\medskip
\ifdefined\hypertarget\hypertarget{sourcepointer-5}{}\fi
\noindent\textit{Source reference.} Original lemma 054F: sources/stacks-chow/tags.txt, source lines 6638--6638. Cited at proof lines 1006-1008: realizes the specialization s' leading to s by a valuation ring map.\par\medskip
\ifdefined\hypertarget\hypertarget{sourcepointer-6}{}\fi
\noindent\textit{Source reference.} Original lemma 0H4K: sources/stacks-chow/tags.txt, source lines 21355--21355. Cited at proof lines 1012-1015: transports the specialization identity along the composite.\par\medskip
\ifdefined\hypertarget\hypertarget{sourcepointer-7}{}\fi
\noindent\textit{Source reference.} Original section 0H4E: sources/stacks-chow/tags.txt, source lines 21349--21349. Cited at proof lines 1017-1018: injectivity of the closed-fibre base change map lets the identity descend.\par\medskip
\begin{lemma}
\label{lemma-relative-cycles-h-descent}
Let $f : X \to S$ be a morphism of schemes. Assume $S$ locally Noetherian
and $f$ locally of finite type. Let $r \geq 0$ be an integer. Let $\alpha$
be a family of $r$-cycles on fibres of $X/S$. Let $\{g_i : S_i \to S\}$
be a h covering (More on Flatness, Definition
\ref{flat-definition-h-covering}). Then $\alpha$ is a relative $r$-cycle
if and only if each base change $g_i^*\alpha$ is a relative $r$-cycle.
\end{lemma}

\begin{proof}
If $\alpha$ is a relative $r$-cycle, then each base change $g_i^*\alpha$ is a
relative $r$-cycle by Lemma \ref{lemma-relative-cycle-functoriality}.
Assume each $g_i^*\alpha$ is a relative $r$-cycle.
Let $g : S' \to S$ be a morphism where $S'$ is the spectrum of a discrete
valuation ring. After replacing $S$ by $S'$, $X$ by $X' = X \times_S S'$, and
$\alpha$ by $\alpha' = g^*\alpha$ and using that the base change of a
h covering is a h covering (More on Flatness, Lemma \ref{flat-lemma-h})
we reduce to the problem studied in the next paragraph.

\medskip\noindent
Assume $S$ is the spectrum of a discrete valuation ring with closed
point $0$ and generic point $\eta$. We have to show that
$sp_{X/S}(\alpha_\eta) = \alpha_0$. Since a h covering is a V covering
(by definition), there is an $i$ and a specialization $s' \leadsto s$
of points of $S_i$ with $g_i(s') = \eta$ and $g_i(s) = 0$, see
Topologies, Lemma \ref{topologies-lemma-refine-qcqs-V}. By
Properties, Lemma \ref{properties-lemma-locally-Noetherian-specialization-dvr}
we can find a morphism $h : S' \to S_i$ from the
spectrum $S'$ of a discrete valuation ring which maps
the generic point $\eta'$ to $s'$ and maps
the closed point $0'$ to $s$. Denote $\alpha' = h^*g_i^*\alpha$.
By assumption we have $sp_{X'/S'}(\alpha'_{\eta'}) = \alpha'_{0'}$.
Since $g = g_i \circ h : S' \to S$ is the morphism of schemes
induced by an extension of discrete valuation rings we conclude that
$sp_{X/S}$ and $sp_{X'/S'}$ are compatible with base change maps
on the fibres, see Lemma \ref{lemma-specialization-extension}.
We conclude that $sp_{X/S}(\alpha_\eta) = \alpha_0$ because
the base change map $Z_r(X_0) \to Z_r(X'_{0'})$ is injective
as discussed in Section \ref{section-relative-fields}.
\end{proof}

\end{document}
